\section{Phase 1 --- Chọn cách lấy đoạn văn}

\subsection{Thiết kế: giải đấu phân tầng}

Lưới đầy đủ của các lựa chọn (retriever $\times$ reranker $\times$ kích thước
chunk) là 72 tổ hợp. Thay vào đó, dự án chạy một giải đấu ba vòng, mỗi vòng chỉ
đưa quán quân đi tiếp --- tổng cộng 23 cấu hình.

\begin{center}\small
\begin{tabularx}{0.94\textwidth}{@{}p{1.4cm} p{5.6cm} L@{}}
\toprule
\textbf{Vòng} & \textbf{Thử gì} & \textbf{Thắng} \\
\midrule
1 & 4 mô hình dense $\times$ 4 mô hình sparse & BGE-M3 learned sparse \\
2 & 3 retriever $\times$ 3 reranker & $+$ \texttt{BAAI/bge-reranker-large} \\
3 & 3 kích thước chunk (256/512/1024) & 512/64 --- \emph{không ai thắng rõ} \\
\bottomrule
\end{tabularx}
\end{center}

Đây là một đánh đổi có ý thức và cần nêu rõ: cách này rẻ hơn nhiều lần, nhưng
\textbf{không phát hiện được tương tác} giữa các lựa chọn ở hai vòng khác nhau.
Ví dụ, khả năng một mô hình dense nào đó phối hợp đặc biệt tốt với chunk 1024
sẽ không bao giờ được kiểm tra, vì dense đã bị loại từ vòng~1.

\subsection{Vòng 1 --- Sparse áp đảo dense}

\begin{center}\small
\begin{tabularx}{0.94\textwidth}{@{}L p{2.2cm} p{3.4cm}@{}}
\toprule
\textbf{So sánh} (nDCG@5, \texttt{resolved}, \NumNDev{} câu) & \textbf{$\Delta$} & \textbf{CI95} \\
\midrule
BGE-M3 sparse \emph{vs} e5-base-v2 dense & \NumSparseVsDense & \NumSparseVsDenseCI \\
BGE-M3 \emph{vs} BM25-okapi-stemmed & \NumBgeMThreeVsBmTwoFive & \NumBgeMThreeVsBmTwoFiveCI \\
e5-base-v2 \emph{vs} bge-small-en-v1.5 & \NumDenseVsDense & \NumDenseVsDenseCI \\
\bottomrule
\end{tabularx}
\end{center}

Chỉ dòng đầu vượt được cả hai điều kiện phán quyết. Khoảng cách
\NumSparseVsDense{} lớn hơn biên độ nhiễu \NumNoiseFloor{} và khoảng tin cậy
không chứa 0. EDA đã giải thích trước vì sao: câu hỏi NewsQA bám vào tên riêng,
ngày tháng và con số --- đúng thứ mà khớp từ vựng bắt được chính xác, còn
embedding dày đặc thì làm nhòe đi thành ``ý nghĩa gần giống''.

Hai dòng sau có khoảng tin cậy \textbf{chứa 0}. Không phân định được BGE-M3 với
BM25 đã stem, và cũng không phân định được hai mô hình dense với nhau.

\keypoint{\textbf{Không tồn tại ``mô hình dense tốt nhất'' trong thí nghiệm này},
và đó tự nó là một kết quả cần báo cáo. Xếp hạng chúng theo số thập phân thứ ba
sẽ là đọc nhiễu thành tín hiệu.}

\subsection{Vòng 2 --- Reranker, và một kết quả phủ định}

Reranker cross-encoder thêm $+0{,}0659$ nDCG@5, lọc top~\NumPoneTopK{} xuống
top~\NumPoneRerankTopN{}. Mức cải thiện này được EDA dự báo trước: với trung vị
\NumDistractorMedian{} đoạn đối thủ cùng chủ đề cho mỗi câu hỏi, việc xếp hạng
lại bằng một mô hình đọc được cả cặp (câu hỏi, đoạn) là bắt buộc.

Hybrid dung hợp dense và sparse \textbf{không chứng minh được là có lợi}. Cần
phát biểu chính xác: kết quả là ``không chứng minh được có lợi'', \emph{không
phải} ``đã chứng minh có hại''. Trong điều kiện dense yếu hơn sparse tới
\NumSparseVsDense{} nDCG@5, việc trộn một tín hiệu yếu vào một tín hiệu mạnh
khó có đường tạo ra lợi ích.

\subsection{Vòng 3 --- Một kết quả null}

\begin{center}\small
\begin{tabular}{@{}l r@{}}
\toprule
Chunk 512 / 64 & \NumChunkFiveOneTwoNdcg \\
Chunk 1024 / 128 & \NumChunkOneZeroTwoFourNdcg \\
Chunk 256 / 32 & \NumChunkTwoFiveSixNdcg \\
\midrule
Khoảng cách cao nhất -- thấp nhất & \NumChunkSpread \\
\bottomrule
\end{tabular}
\end{center}

\NumChunkSpread{} nằm dưới biên độ nhiễu \NumNoiseFloor{}, và cả ba khoảng tin
cậy chồng lấn nhau. Không có cơ sở thống kê để xếp hạng ba cấu hình này.

Lý do nằm ở phân bố độ dài bài báo: sau khi phục hồi, trung vị vẫn là 724 token
mỗi bài và \textbf{75{,}9\% số bài vẫn chỉ có tối đa 2 chunk}. Với ba phần tư
kho ngữ liệu, bài toán tìm đúng chunk gần như đồng nhất với bài toán tìm đúng
bài báo, nên dư địa để chiến lược chunking tạo khác biệt là rất hẹp. Việc chọn
512/64 dựa trên lý do vận hành --- không quá ngắn khiến câu trả lời bị cắt đôi,
không quá dài gây loãng và tốn token ở Phase~2 --- chứ không phải vì nó thắng.

\subsection{Ablation bổ sung --- Contextual chunking}

Vòng 3 kết luận kích thước chunk không quan trọng. Câu hỏi tiếp theo: nếu không
phải cắt to hay cắt nhỏ, thì \emph{viết lại nội dung} chunk có quan trọng
không? Ablation thêm \NumCtxChars{} ký tự ngữ cảnh của bài báo mẹ vào đầu mỗi
chunk (\NumCtxContextualised{}/\NumCorpusChunks{} chunk được thêm; chunk đầu
tiên bỏ qua vì đã tự mang ngữ cảnh).

\begin{center}\small
\begin{tabularx}{0.94\textwidth}{@{}L p{2cm} p{2cm} p{3.4cm}@{}}
\toprule
nDCG@5 (\texttt{resolved}, \NumCtxNSamples{} câu) & \textbf{plain} & \textbf{contextual} & \textbf{$\Delta$ và CI95} \\
\midrule
Dense (e5-base-v2) & \NumCtxDensePlain & \NumCtxDenseCtx & \NumCtxEffectDense{} \NumCtxEffectDenseCI \\
Sparse (BGE-M3) & \NumCtxSparsePlain & \NumCtxSparseCtx & \NumCtxEffectSparse{} \NumCtxEffectSparseCI \\
\bottomrule
\end{tabularx}
\end{center}

Hai hiệu ứng ngược chiều, cả hai đều có ý nghĩa thống kê. Contextual chunking
\textbf{giúp dense} và \textbf{hại sparse} --- và cả hai đều nhất quán với phát
hiện của EDA. Với dense, \NumCtxChars{} ký tự ngữ cảnh là thông tin embedding
trước đó không có: nó giúp phân biệt hai đoạn gần giống nhau đến từ hai bài
khác nhau. Với sparse, chính \NumCtxChars{} ký tự đó lặp lại trên mọi chunk của
cùng một bài, \textbf{làm loãng tần suất của những từ hiếm} vốn là thứ giúp
BGE-M3 tìm đúng đoạn.

Kết quả này \textbf{không đổi quyết định của Phase~1}: sparse vẫn thắng dense
trên cả hai kho, khoảng cách thu hẹp từ \NumCtxGapPlain{} xuống
\NumCtxGapCtx{} nhưng không đảo ngược. Ngoài ra, cả \NumCtxEffectDense{} lẫn
\NumCtxEffectSparse{} đều nhỏ hơn ngưỡng thực dụng \NumNoiseFloor{}.

Hai điểm cần ghi nhận trung thực. Thứ nhất, đây là ablation chứ không phải một
vòng của giải đấu: nó chạy trên harness riêng (\texttt{top\_k} 10, không
reranker, dense dùng tích ma trận chính xác thay vì HNSW nên tái lập được), nên
con số tuyệt đối không so trực tiếp với bảng vòng~1. Thứ hai, ablation chạy
trên cả \NumCtxNSamples{} câu, tức \textbf{có chạm vào các câu held-out của
Phase~2}; vì đây là chỉ số truy xuất thuần và không có quyết định cấu hình nào
được rút ra từ nó, phần sinh trên held-out vẫn chưa bị nhìn thấy.

\subsection{Cấu hình khóa}

\begin{center}\small
BGE-M3 learned sparse (\texttt{top\_k=\NumPoneTopK}) $\rightarrow$
\texttt{BAAI/bge-reranker-large} (top \NumPoneRerankTopN{}) $\rightarrow$ chunk 512/64
\end{center}

\begin{center}\small
\begin{tabular}{@{}c c c c c@{}}
\toprule
\textbf{Hit@1} & \textbf{Hit@3} & \textbf{Hit@5} & \textbf{nDCG@5} & \textbf{Latency P50} \\
\midrule
\NumPoneHitOne & \NumPoneHitThree & \NumPoneHitFive & \NumPoneNdcgFive & \NumPoneLatency{} ms \\
\bottomrule
\end{tabular}
\end{center}

Chỉ mục được đóng gói và băm (\texttt{\NumLockedIndexSha\ldots}). Con số cần
giữ lại cho phần sau: \textbf{Hit@5 $=$ \NumPoneHitFive{}}. Nó sẽ trở thành
trần của Phase~2.
